%%% Внесите свои данные - Input your data
%%
%%
\newcommand{\Author}{Р.В.\,Губин} % И.О. Фамилия автора 
\newcommand{\AuthorFull}{Губин Роман Витальевич} % Фамилия Имя Отчество автора
\newcommand{\AuthorFullDat}{Губину Роману Витальевичу} % Фамилия Имя Отчество автора в дательном падеже (Кому? Студенту...)
\newcommand{\AuthorFullVin}{Губина Романа Витальевича} % в винительном падеже (Кого? что?  Програмиста ...)
\newcommand{\AuthorPhone}{+7-992-202-60-24} % номер телефорна автора для оперативной связи  
\newcommand{\Supervisor}{Н.И.\,Зайцева} % И. О. Фамилия научного руководителя
\newcommand{\SupervisorFull}{Зайцева Надежда Игоревна} % Фамилия Имя Отчество научного руководителя
\newcommand{\SupervisorVin}{Н.И.\,Зайцеву} % И. О. Фамилия научного руководителя  в винительном падеже (Кого? что? Руководителя ...)
\newcommand{\SupervisorJob}{доцент} %
\newcommand{\SupervisorJobVin}{доцента} % в винительном падеже (Кого? что?  Програмиста ...)
\newcommand{\SupervisorDegree}{кандидат технических наук} %
\newcommand{\SupervisorTitle}{} % 
%%
%%
%Руководитель, утверждающий задание
\newcommand{\Head}{К.Н.\,Козлов} % И. О. Фамилия руководителя подразделения (руководителя ОП)
\newcommand{\HeadDegree}{Руководитель ОП}% Только должность:   
%Руководитель ОП % <- ПРАВИЛЬНЫЙ ВАРИАНТ ДЛЯ 09.03.03 и 09.04.03
%Заведующий % кафедрой
%Директор % Высшей школы
%Зам. директора
\newcommand{\HeadDep}{} % заменить на краткую аббревиатуру подразделения или ОСТАВИТЬ ПУСТЫМ, если утверждает руководитель ОП

%%% Руководитель, принимающий заявление
\newcommand{\HeadAp}{И.О.\,Фамилия} % И. О. Фамилия руководителя подразделения (руководителя ОП)
\newcommand{\HeadApDegree}{Должность руководителя}% Только должность:   
%Руководитель ОП 
%Заведующий кафедрой
%Директор Высшей школы
\newcommand{\HeadApDep}{} % заменить на краткую аббревиатуру подразделения или оставить пустым, если утверждает руководитель ОП
%%% Консультант по нормоконтролю
\newcommand{\ConsultantNorm}{Л.А.\,Арефьева} % И. О. Фамилия консультанта по нормоконтролю. ТОЛЬКО из числа ППС!
\newcommand{\ConsultantNormDegree}{должность, степень} %   
%%% Первый консультант
\newcommand{\ConsultantExtraFull}{Фамилия Имя Отчетство} % Фамилия Имя Отчетство дополнительного консультанта 
\newcommand{\ConsultantExtra}{И.О.\,Фамилия} % И. О. Фамилия дополнительного консультанта 
\newcommand{\ConsultantExtraDegree}{должность, степень} % 
\newcommand{\ConsultantExtraVin}{И.О.\,Фамилию} % И. О. Фамилия дополнительного консультанта в винительном падеже (Кого? что? Руководителя ...)
\newcommand{\ConsultantExtraDegreeVin}{должность, степень} %  в винительном падеже (Кого? что? Руководителя ...)
%%% Второй консультант
\newcommand{\ConsultantExtraTwoFull}{Фамилия Имя Отчетство} % Фамилия Имя Отчетство дополнительного консультанта 
\newcommand{\ConsultantExtraTwo}{И.О.\,Фамилия} % И. О. Фамилия дополнительного консультанта 
\newcommand{\ConsultantExtraTwoDegree}{должность, степень} % 
\newcommand{\ConsultantExtraTwoVin}{И.О.\,Фамилию} % И. О. Фамилия дополнительного консультанта в винительном падеже (Кого? что? Руководителя ...)
\newcommand{\ConsultantExtraTwoDegreeVin}{должность, степень} %  в винительном падеже (Кого? что? Руководителя ...)
\newcommand{\Reviewer}{И.О.\,Фамилия} % И. О. Фамилия резензента. Обязателен только для магистров.
\newcommand{\ReviewerDegree}{должность, степень} % 
%%
%%
\renewcommand{\thesisTitle}{Разработка алгоритма функционирования самоорганизующейся низкоскоростной системы передачи данных}
\newcommand{\thesisDegree}{работа бакалавра}% дипломный проект, дипломная работа, магистерская диссертация %c 2020
\newcommand{\thesisTitleEn}{Development of an Operating Algorithm for Self-Organizing Low-Rate Data Transmission System} %2020
\newcommand{\thesisDeadline}{дд.мм.202X} %Последний день преддипломной практики согласно учебному плану.
\newcommand{\thesisStartDate}{дд.мм.202X}
\newcommand{\thesisYear}{2025} % заменить на год защиты
\newcommand{\approveYear}{202\underline{\hspace*{0.01\textheight}}} % <- НЕ МЕНЯТЬ, ТОЛЬКО С 2030го :)
%%
%%
\newcommand{\group}{5030102/10401} % заменить вместо N номер группы
\newcommand{\thesisSpecialtyCode}{01.03.02}% код направления подготовки
\newcommand{\thesisSpecialtyTitle}{Прикладная математика и информатика} % наименование направления/специальности
\newcommand{\thesisOPPostfix}{04} % последние цифры кода образовательной программы (после <<_>>)
\newcommand{\thesisOPTitle}{Биоинформатика}% наименование образовательной программы
%%
%%
\newcommand{\institute}{
Физико-механический институт
%Институт компьютерных наук и~кибербезопасности
%Гуманитарный институт
%Инженерно-строительный институт
%Институт биомедицинских систем и технологий
%Институт металлургии, машиностроения и транспорта
%Институт передовых производственных технологий
%Институт прикладной математики и механики
%Институт физики, нанотехнологий и телекоммуникаций
%Институт физической культуры, спорта и туризма
%Институт энергетики и транспортных систем
%Институт промышленного менеджмента, экономики и торговли
}%
%%
%%




%%% Задание ключевых слов и аннотации
%%
%%
%% Ключевых слов от 3 до 5 слов или словосочетаний в именительном падеже именительном падеже множественного числа (или в единственном числе, если нет другой формы) по правилам русского языка!!!
%%
%%
\newcommand{\keywordsRu}{%
самоорганизующиеся сети, %
низкоскоростная передача данных, %
P2P-сети, %
абоненты, %
алгоритмы маршрутизации, %
радиоканал, %
синхронизация времени%
}%
%%
%%
\newcommand{\keywordsEn}{%
self-organizing networks, %
low-rate data transmission, %
peer-to-peer networks, %
network nodes, %
routing algorithms, %
wireless channel, %
time synchronization%
}%
%%
%%
%% Реферат ОТ 1000 ДО 1500 знаков на русский или английский текст
%%
%Реферат должен содержать:
%- предмет, тему, цель ВКР;
%- метод или методологию проведения ВКР:
%- результаты ВКР:
%- область применения результатов ВКР;
%- выводы.
\newcommand{\abstractRu}{%
В работе исследуется алгоритм функционирования самоорганизующейся низкоскоростной системы передачи данных для мобильных одноранговых сетей с радиоканалом дальностью до 10 км и скоростью передачи 32 бит/с – 37 кбит/с.%

\textbf{Цель работы} — разработка устойчивого алгоритма, обеспечивающего:  \\
- динамическое обнаружение абонентов при скорости их перемещения до 60 км/ч,  \\
- минимизацию потерь данных за счёт адаптивных повторных передач,  \\
- синхронизацию времени с погрешностью <1 мс.  

\textbf{Методология} включает:  \\
- Моделирование движения абонентов,  \\
- Оптимизацию формата передаваемых данных,  \\
- Разработка алгоритма маршрутизации.  

\textbf{Результаты}:  \\
- Для сети из 5 абонентов достигнута задержка передачи 1.5 сек при 0\% потерь,  \\
- При 15 абонентах и скорости 60 км/ч время синхронизации составило 20 сек (потери фреймов – 5\%),  \\
- Максимальное число абонентов: 20 при 37 кбит/с, 8-12 при 32 бит/с.  

\textbf{Область применения}: военные сети, интернет вещей (\textit{англ.} Internet of Things, \textit{далее} — IoT) в удалённых районах, автономные транспортные системы.  

\textbf{Выводы}: предложенный алгоритм обеспечивает надёжность сети в условиях низкой скорости передачи и мобильности абонентов, то есть её устойчивость к помехам и динамическом изменении положении абонентов.
}
%%
%%
\newcommand{\abstractEn}{%
This research investigates an algorithm for self-organizing low-speed data transmission systems in mobile peer-to-peer (P2P) networks with a radio channel range of up to 10 km and data rates of 32 bps to 37 kbps.

\textbf{Research objectives} include the development of a robust algorithm that ensures:
\begin{enumerate}
	\item Dynamic node discovery at movement speeds up to 60 km/h
	\item Data loss minimization through adaptive retransmission
	\item Time synchronization with error <1 ms
\end{enumerate}

\textbf{Methodology} comprises:
\begin{itemize}
	\item Node movement modeling (coordinates $x(t)$, $y(t)$ and communication radius $R \leq 10$ km)
	\item Frame format optimization (16-bit header + 128-bit payload + CRC)
	\item Hybrid routing with reachability zone prediction
\end{itemize}

\textbf{Results}:
\begin{itemize}
	\item For a 8-node network: achieved 1.5 sec transmission delay with 0\% packet loss
	\item For 15 nodes at 60 km/h: synchronization time of 20 sec (frame loss rate 5\%)
	\item Maximum network capacity: 20 nodes at 37 kbps, 8-12 nodes at 32 bps
\end{itemize}

\textbf{Applications}: military networks, IoT in remote areas, autonomous transportation systems.

\textbf{Conclusions}: The proposed algorithm ensures network survivability under low-speed transmission and node mobility conditions, but requires clustering for scalability.
}


%%% РАЗДЕЛ ДЛЯ ОФОРМЛЕНИЯ ПРАКТИКИ
%Место прохождения практики
\newcommand{\PracticeType}{Отчет о прохождении % 
	%стационарной производственной (технологической (проектно-технологической)) %
	преддипломной % тип и вид ЗАМЕНИТЬ
	практики}

\newcommand{\Workplace}{СПбПУ, ФизМех, ВШ ПМ и ВФ} % TODO Rename this variable

% Даты начала/окончания
\newcommand{\PracticeStartDate}{%
02.02.2025%
%	22.06.2020
}%
\newcommand{\PracticeEndDate}{%
	07.05.2025%
%	18.07.2020%
}%
%%

\newcommand{\School}{
	Высшая школа прикладной математики и вычислительной физики 
%	Высшая школа программной инженерии 
}
\newcommand{\practiceTitle}{\thesisTitle}


%% ВНИМАНИЕ! Необходимо либо заменить текст аннотации (ключевых слов) на русском и английском, либо удалить там весь текст, иначе в свойства pdf-отчета по практике пойдет шаблонный текст.

%%% Не меняем дальнейшую часть - Do not modify the rest part
%%
%%
%%
%%
\ifnumequal{\value{docType}}{1}{% Если ВКР, то...
	\newcommand{\DocType}{Выпускная квалификационная работа}
	\newcommand{\pdfDocType}{\DocType~(\thesisDegree)} %задаём метаданные pdf файла
	\newcommand{\pdfTitle}{\thesisTitle}
}{% Иначе 
	\newcommand{\DocType}{\PracticeType}
	\newcommand{\pdfDocType}{\DocType} %задаём метаданные pdf файла
	\newcommand{\pdfTitle}{\practiceTitle}
}%
\newcommand{\HeadTitle}{\HeadDegree~\HeadDep}
\newcommand{\HeadApTitle}{\HeadApDegree~\HeadApDep}
\newcommand{\thesisOPCode}{\thesisSpecialtyCode\_\thesisOPPostfix}% код образовательной программы
\newcommand{\thesisSpecialtyCodeAndTitle}{\thesisSpecialtyCode~\thesisSpecialtyTitle}% Код и наименование направления/специальности
\newcommand{\thesisOPCodeAndTitle}{\thesisOPCode~\thesisOPTitle} % код и наименование образовательной программы
%%
%%
\hypersetup{%часть болка hypesetup в style
		pdftitle={\pdfTitle},    % Заголовок pdf-файла
		pdfauthor={\AuthorFull},    % Автор
		pdfsubject={\pdfDocType. Шифр и наименование направления подготовки: \thesisSpecialtyCodeAndTitle.},      % Тема
		pdfcreator={LaTeX, SPbPU-student-thesis-template},     % Приложение-создатель
%		pdfproducer={},  % Производитель, Производитель PDF % будет выставлена автоматически
		pdfkeywords={\keywordsRu}
}
%%
%%
%% вспомогательные команды
\newcommand{\firef}[1]{рис.\ref{#1}} %figure reference
\newcommand{\taref}[1]{табл.\ref{#1}}	%table reference
%%
%%
%% Архивный вариант задания ключевых слов, аннотации и благодарностей 
% Too hard to export data from the environment to pdf-info
% https://tex.stackexchange.com/questions/184503/collecting-contents-of-environment-and-store-them-for-later-retrieval
%заменить NewEnviron на newenvironment для распознавания команды в TexStudio
%\NewEnviron{keywordsRu}{\noindent\MakeUppercase{\BODY}}
%\NewEnviron{keywordsEn}{\noindent\MakeUppercase{\BODY}}
%\newenvironment{abstractRu}{}{}
%\newenvironment{abstractEn}{}{}
%\newenvironment{acknowledgementsRu}{\par{\normalfont \acknowledgements.}}{}
%\newenvironment{acknowledgementsEn}{\par{\normalfont \acknowledgementsENG.}}{}


%%% Переопределение именований %%% Не меняем - Do not modify
%\newcommand{\Ministry}{Минобрнауки России} 
\newcommand{\Ministry}{Министерство науки и высшего образования Российской~Федерации} %с 2020
\newcommand{\SPbPU}{Санкт-Петербургский политехнический университет Петра~Великого}
\newcommand{\SPbPUOfficialPrefix}{Федеральное государственное автономное образовательное учреждение высшего образования}
\newcommand{\SPbPUOfficialShort}{ФГАОУ~ВО~<<СПбПУ>>}
%% Пробел между И. О. не допускается.
\renewcommand{\alsoname}{см. также}
\renewcommand{\seename}{см.}
\renewcommand{\headtoname}{вх.}
\renewcommand{\ccname}{исх.}
\renewcommand{\enclname}{вкл.}
\renewcommand{\pagename}{Pages}
\renewcommand{\partname}{Часть}
\renewcommand{\abstractname}{\textbf{Аннотация}}
\newcommand{\abstractnameENG}{\textbf{Annotation}}
\newcommand{\keywords}{\textbf{Ключевые слова}}
\newcommand{\keywordsENG}{\textbf{Keywords}}
\newcommand{\acknowledgements}{\textbf{Благодарности}}
\newcommand{\acknowledgementsENG}{\textbf{Acknowledgements}}
\renewcommand{\contentsname}{Content} % 
%\renewcommand{\contentsname}{Содержание} % (ГОСТ Р 7.0.11-2011, 4)
%\renewcommand{\contentsname}{Оглавление} % (ГОСТ Р 7.0.11-2011, 4)
\renewcommand{\figurename}{Рис.} % Стиль СПбПУ
%\renewcommand{\figurename}{Рисунок} % (ГОСТ Р 7.0.11-2011, 5.3.9)
\renewcommand{\tablename}{Таблица} % (ГОСТ Р 7.0.11-2011, 5.3.10)
%\renewcommand{\indexname}{Предметный указатель}
\renewcommand{\listfigurename}{Список рисунков}
\renewcommand{\listtablename}{Список таблиц}
\renewcommand{\refname}{\fullbibtitle}
\renewcommand{\bibname}{\fullbibtitle}

\newcommand{\chapterEnTitle}{Сhapter title} % <- input the English title here (only once!) 
\newcommand{\chapterRuTitle}{Название главы}          % <- введите 
\newcommand{\sectionEnTitle}{Section title} %<- input subparagraph title in english
\newcommand{\sectionRuTitle}{Название подраздела} % <- введите название подраздела по-русски
\newcommand{\subsectionEnTitle}{Subsection title} % - input subsection title in english
\newcommand{\subsectionRuTitle}{Название параграфа} % <- введите название параграфа по-русски
\newcommand{\subsubsectionEnTitle}{Subsubsection title} % <- input subparagraph title in english
\newcommand{\subsubsectionRuTitle}{Название подпараграфа} % <- введите название подпараграфа по-русски